% Copyright (c) 2026 Twin Vector Labs LLC.
% All rights reserved.
%
% Research-direction companion to spin_readout.tex.  Reframes the
% composite-spin readout through the Cartan (KAK) decomposition: the
% local--nonlocal--local structure, the identification of the SU(3)
% color Weyl group with the S_3 of the J^2 swaps, and the question of
% whether the entangling exponential plays the role of a gluon.
\documentclass[11pt,oneside]{article}

\usepackage[margin=1in]{geometry}
\usepackage{amsmath}
\usepackage{amssymb}
\usepackage{amsthm}
\usepackage{mathtools}
\usepackage{braket}
\usepackage{bm}
\usepackage{enumitem}
\usepackage{booktabs}
\usepackage{array}
\usepackage{xcolor}
\usepackage[colorlinks=true,linkcolor=blue!55!black,citecolor=green!45!black,%
            urlcolor=blue!55!black]{hyperref}

% --- convenience macros (mirrors spin_readout.tex) ---------------------------
\newcommand{\R}{\mathbb{R}}
\newcommand{\C}{\mathbb{C}}
\newcommand{\Z}{\mathbb{Z}}
\newcommand{\halff}{\tfrac{1}{2}}
\newcommand{\id}{\mathbb{1}}
\newcommand{\Sig}{\Sigma}
\newcommand{\Tr}{\operatorname{Tr}}
\newcommand{\eqdef}{\vcentcolon=}
\DeclareMathOperator{\sgn}{sgn}
\newcommand{\expval}[1]{\langle #1 \rangle}
\newcommand{\mel}[3]{\langle #1 | #2 | #3 \rangle}
\newcommand{\norm}[1]{\left\lVert #1 \right\rVert}
\newcommand{\code}[1]{\texttt{#1}}
\newcommand{\up}{\uparrow}
\newcommand{\down}{\downarrow}
% additions for this companion
\newcommand{\SU}{\mathrm{SU}}
\newcommand{\SO}{\mathrm{SO}}
\newcommand{\Spin}{\mathrm{Spin}}
\newcommand{\Weyl}{\mathcal{W}}
\newcommand{\fk}{\mathfrak{k}}
\newcommand{\fa}{\mathfrak{a}}
\newcommand{\fp}{\mathfrak{p}}
\newcommand{\fg}{\mathfrak{g}}
\newcommand{\Pin}{P_{\mathrm{in}}}
\newcommand{\Pout}{P_{\mathrm{out}}}
\theoremstyle{plain}
\newtheorem{claim}{Claim}
\newtheorem{remark}{Remark}

\title{\bfseries A Cartan / Weyl--Chamber Reading of the Composite--Spin Readout\\[4pt]
\large Local--nonlocal--local transport between color holes, and whether the
entangling exponential is a gluon}
\author{Tessera theory notes \quad(companion to \code{spin\_readout.tex}; part of issue \#410)}
\date{\today}

\begin{document}
\maketitle

\begin{abstract}
\noindent
This is a research--direction companion to \code{spin\_readout.tex} (same
directory), which establishes that the proton's composite spin $J^2=\tfrac34$ is
an \emph{entangled}, mixed--symmetry three--body quantity and that any readout
reducing each color hole to its own Bloch vector provably floors at the product
mixture ($\tfrac74$ per ket, ${\sim}\tfrac94$ for three). Here we take seriously a
single organizing idea: read the relation between two color holes through the
\emph{Cartan} (KAK) decomposition $g = k_1\,a\,k_2$, whose structure is exactly
\emph{local--nonlocal--local}. Three things fall out. (i) The product floor is
precisely the statement that the reconstructed joint state is confined to the
$K$--orbit (the $k$'s), discarding the abelian/entangling middle factor $a$. (ii)
The Weyl group of the color group $\SU(3)$ is the symmetric group $S_3$, and it is
the \emph{same} $S_3$ whose transpositions appear as the swap operators
$P_{12},P_{13},P_{23}$ in $J^2=\tfrac34+\sum P_{ij}$; the carried color register
$[1,\omega,\omega^2]$ is the Weyl--covariant assignment of weights to holes. (iii)
The entangling exponential $\exp(i\sum_a c_a\,\sigma_a\!\otimes\!\sigma_a)$ already
appears in the codebase --- in the abelian (Schwinger/QED) interaction model,
where it is explicitly the \emph{neutral photon}. ``Acts like a gluon'' is then a
precise, falsifiable proposal: occupy the same mediator slot with the
\emph{non--abelian}, color--valued, geometry--sourced version. We state what would
make that true, what is merely analogy, and a concrete next experiment. Nothing
here is claimed as established; it is a lens and a program.
\end{abstract}

\tableofcontents
\bigskip

% =============================================================================
\section{What this companion adds, in one paragraph}
\label{sec:premise}
% =============================================================================

\code{spin\_readout.tex} ends at the \emph{total--space spin operator}: to recover
$\tfrac34$ the spin Casimir must act on the whole carried representative at once,
keeping the cross--hole terms. That is correct, and we do not retract it. What
this note adds is a \emph{name} for the missing object and a \emph{coordinate} for
it. The name is the abelian middle factor $a$ of a Cartan decomposition; the
coordinate is its Weyl--chamber position (equivalently, the local--unitary
invariants of the two--hole state). Naming it buys three concrete things: a
sharper diagnosis of why the product read floors (\S\ref{sec:kak}), an identity
--- not an analogy --- between the color group's Weyl symmetry and the swap
operators that produce $\tfrac34$ (\S\ref{sec:weyl}), and a precise version of the
user's question ``does the exponential act like a gluon?'' (\S\ref{sec:gluon}).
Two payoffs (\S\ref{sec:payoff}) and the honest caveats (\S\ref{sec:caveats})
follow.

% =============================================================================
\section{The Cartan (KAK) decomposition and the entangling core}
\label{sec:kak}
% =============================================================================

\paragraph{The decomposition.}
Let $G$ be a compact Lie group with Cartan involution $\theta$, splitting the
algebra as $\fg=\fk\oplus\fp$ ($[\fk,\fk]\subseteq\fk$, $[\fk,\fp]\subseteq\fp$,
$[\fp,\fp]\subseteq\fk$). Choose a maximal abelian $\fa\subseteq\fp$. The
\emph{Cartan decomposition} (the $KAK$ decomposition) writes every $g\in G$ as
\begin{equation}
  g = k_1\,a\,k_2, \qquad k_1,k_2\in K=\exp\fk,\quad a\in A=\exp\fa .
  \label{eq:kak}
\end{equation}
The $A$ factor is determined up to the \emph{Weyl group} $\Weyl=N_K(\fa)/Z_K(\fa)$,
and a fundamental domain for $\Weyl$ acting on $\fa$ is a \emph{Weyl chamber}. The
chamber coordinates of $a$ are the only data of $g$ invariant under
$g\mapsto k_1' g\,k_2'$ --- the genuinely two--sided, frame--independent content.

\paragraph{The two--qubit instance (the one we use).}
For two spin--$\halff$ holes, $G=\SU(4)$, $K=\SU(2)\times\SU(2)$, and
\eqref{eq:kak} is the standard ``canonical'' / KAK form of a two--qubit gate
\cite{khaneja,kraus}:
\begin{equation}
  \boxed{\;U = \underbrace{(k_1\otimes k_2)}_{\text{local}}\;
              \underbrace{\exp\!\Big(i\textstyle\sum_{a=1}^{3} c_a\,
                          \sigma_a\!\otimes\!\sigma_a\Big)}_{\text{nonlocal }=\,A}\;
              \underbrace{(k_3\otimes k_4)}_{\text{local}}\;}
  \label{eq:kak2}
\end{equation}
with $k_i\in\SU(2)$ acting on a single hole and $(c_1,c_2,c_3)$ the Weyl--chamber
point ($\tfrac{\pi}{4}\ge c_1\ge c_2\ge|c_3|$). The structure is exactly the
\emph{local--nonlocal--local} pattern of the user's intuition. The load--bearing
theorem:
\begin{claim}[Only $A$ can change entanglement]\label{cl:onlyA}
Local unitaries $k\otimes k'$ preserve the entanglement (Schmidt spectrum) of any
two--hole state. Hence the entanglement created by $U$ depends on $U$ only through
its chamber point $(c_1,c_2,c_3)$, i.e.\ through $A$ alone.
\end{claim}

\paragraph{The product floor, in chamber language.}
\code{spin\_readout.tex} (and \code{joint\_proton\_spin\_findings.md}, \#489)
establish that the $J^2$ \emph{operator} is correct --- fed clean states by hand
it returns $\tfrac34,\ \tfrac74,\ \tfrac{15}{4}$ (regression guard in
\code{test\_dk\_joint\_spin.py}). The floor is therefore on the \emph{state} side:
\code{composite\_j2} reconstructs the joint spin state as
\begin{equation}
  \ket{\Psi}_{\text{read}} = \ket{s_1}\otimes\ket{s_2}\otimes\ket{s_3},
  \qquad \ket{s_h}=\code{emergent\_spinor}(h),
\end{equation}
a literal \code{np.kron} of three \emph{independently} extracted per--hole
spinors. By Claim~\ref{cl:onlyA} a product is fixed at the trivial chamber corner
$c=(0,0,0)$ (the $K$--orbit of a product), and no local post--processing moves it
off; the entangled $J=\halff$ eigenstate lives in a different orbit. So
$\expval{J^2}_{\text{read}}$ is pinned in the three--spin--$\halff$ \emph{product}
range $[\tfrac32,\tfrac{15}{4}]$, clustering near $\tfrac94$ --- exactly the
observed floor. \textbf{The missing ingredient is the $A$ factor in the state
reconstruction}: the inter--hole correlations the build put in and the per--hole
extraction takes out. (This is the same conclusion as the total--space operator of
\code{spin\_readout.tex}; the chamber language just makes the discarded object a
named coordinate.)

\begin{remark}
The single swap $P_{ij}$ is itself \emph{maximally nonlocal} --- it sits at the
far chamber corner $c=(\tfrac\pi4,\tfrac\pi4,\tfrac\pi4)$ (three CNOTs to
synthesize). Note carefully: a swap is maximally nonlocal yet has zero
\emph{entangling power} (it maps products to products). This is consistent with
the fact that $\expval{J^2}$ on the product $\ket{uud}$ is the definite number
$\tfrac74$: evaluating $J^2$ on a product needs no entanglement; reaching the
$\tfrac34$ \emph{eigenstate} does. The chamber distinguishes the two notions
cleanly, which is one reason the language is worth importing.
\end{remark}

\begin{remark}[Two distinct ``Weyl''s; the chiral one carries the complex deficit]
``Weyl'' is used in two unrelated senses, worth separating before
\S\ref{sec:weyl} leans on the first. The Weyl \emph{group}/chamber is the $S_n$ of
the entangling coordinate above (and the $S_3$ of the color register). Distinct
from it is the Weyl \emph{chirality} split of the \emph{single--fiber} frame
rotation: in the Euclidean per--cell embedding $\Spin(4)=\SU(2)_L\times\SU(2)_R$,
and with $P_\pm=\halff(\id\pm\gamma_5)$ the structural spin generators are
block--diagonal in chirality. The \emph{physical} spatial spin is the
\emph{diagonal} $\SU(2)$ --- the same $\halff\sigma_a$ on both blocks --- so
writing a connector's exponent as $\vec c=\vec\theta\pm i\vec\beta$, the real part
$\vec\theta$ is the rotation (spin) and the imaginary part $\vec\beta$ is the
Lorentzian deficit/boost (the $\mathrm{Im}$ of $\vec c$ is the same physics as the
$\mathrm{Im}$ of the complex Regge action). This is entirely on the $K$--type
spin--fiber side: by Claim~\ref{cl:onlyA} it does \emph{not} entangle. Its role is
definitional --- how the per--hole spin generator is fixed consistently with the
complex action --- leaving the entangling content where
\S\ref{sec:weyl}--\ref{sec:gluon} put it: the color $A$ factor, not this chiral
split.
\end{remark}

% =============================================================================
\section{The Weyl group of $\SU(3)_{\text{color}}$ \emph{is} the $S_3$ of the
$J^2$ swaps}
\label{sec:weyl}
% =============================================================================

This section is the one place the Cartan picture is more than analogy. The claim
is an equality of groups, with one physical input (Pauli antisymmetry).

\paragraph{Three standard facts.}
\begin{enumerate}[leftmargin=2em]
  \item The Weyl group of $\SU(n)$ (root system $A_{n-1}$) is the symmetric group
        $S_n$. For color, $\Weyl(\SU(3))=S_3$, order $6$.
  \item The fundamental rep $\mathbf{3}$ has three weights $\mu_1,\mu_2,\mu_3$ at
        the corners of a triangle in the weight plane, with $\sum_i\mu_i=0$;
        $\Weyl=S_3$ permutes them. The three color holes carry one weight each.
  \item The swap form of the spin Casimir (Eq.~(j2swap) of
        \code{spin\_readout.tex}),
        \begin{equation}
          J^2 = \tfrac34\,\id + \big(P_{12}+P_{13}+P_{23}\big),
          \label{eq:j2swap}
        \end{equation}
        is built from the three transpositions of $S_3$ --- its Weyl reflections.
\end{enumerate}

\paragraph{The carried register is the Weyl--covariant assignment.}
The codebase already realizes fact (2): the carried representative
$[1,\omega,\omega^2]$ ($\omega=e^{2\pi i/3}$,
\code{EigenstateSynthesis(st,3).carriedRepresentative}) is \emph{not}
$[1,\omega,\omega^2]$ placed by hand but the $\omega$--eigenvector of the
window--cycling generator $g$, the $\Z_3\subset S_3$ rotation of the color
weights (\code{proton\_symmetric\_windows.tex}, \#398):
\begin{equation}
  \Pin\,\phi = \omega\,\phi, \qquad
  M\,\Pin = \Pout\,M ,
  \label{eq:intertwine}
\end{equation}
where $M$ is the inter--hole transport and $\Pin,\Pout$ are the (signed) color
$3$--cycles on the input/output hole triples. Equation~\eqref{eq:intertwine} says
$M$ \emph{intertwines} the color $\Z_3$ --- exactly the statement that $M$ carries
the Weyl rotation between holes.

\paragraph{The locking (the one physical input).}
On the \emph{physical} baryon the position--permutation $S_3$ and the
color--Weyl $S_3$ are not two groups but one. Pauli antisymmetry forces color into
the antisymmetric singlet $\epsilon_{abc}$ (Eq.~(colorsinglet) of
\code{spin\_readout.tex}); antisymmetry under exchanging the \emph{positions} of
two quarks therefore acts on their \emph{colors} as the corresponding weight
transposition. Hence:
\begin{equation}
  \boxed{\;\text{the swaps }P_{ij}\text{ in }\eqref{eq:j2swap}
          \text{ are the Weyl reflections of }\SU(3)_{\text{color}}
          \text{ acting on the carried register.}\;}
\end{equation}
The $2-1-1=0$ cancellation that yields $\tfrac34$ (\code{spin\_readout.tex},
\S3) is the $S_3$--Weyl structure of the color register read in the spin sector.
This is why ``a Weyl/Cartan move from one hole to the next'' is the right
language: the off--diagonal root generators $E_{\pm\alpha}$ of $\SU(3)$ --- the
color--raising/lowering directions, the genuinely non--abelian gluon directions
--- are precisely what cycle a hole's color $r\!\to\!g\!\to\!b$, and $M$ in
\eqref{eq:intertwine} is the existing discrete seed of that action.

% =============================================================================
\section{Does the exponential act like a gluon?}
\label{sec:gluon}
% =============================================================================

Short answer: it occupies the right \emph{slot}, but in the codebase that slot is
currently filled by the \emph{photon}, and the upgrade to a gluon is a specific,
non--trivial list.

\paragraph{The photon is already there.}
The KAK decomposition \eqref{eq:kak2} is implemented in the quantum/interaction
subsystem (\code{src/simulations/InteractionSimulation.cpp}). The comment at the
construction is explicit:
\begin{quote}\small
\emph{Cartan/local--frame model: under the KAK decomposition
$U=(K_1\!\otimes\!K_2)\cdot\exp(i\,c\,\sigma\sigma)\cdot(K_3\!\otimes\!K_4)$,
\dots\ $\Sigma_{AB}\equiv\exp(i\,c\,\sigma\sigma)$ is the entangling core.}
\end{quote}
and that entangling core is identified physically as ``\emph{the entangling--core
analog of a neutral photon}'' with charge $q_{AB}=0$. So the local--nonlocal--local
exponential is \emph{already} playing the role of a force mediator in the engine
--- as the abelian, neutral one. The user's question is whether the same slot, in
the color sector, is a gluon. That is the right generalization to ask, and it is
sharp because the abelian case is concrete.

\paragraph{What would make it a gluon (and not a photon).}
Three conditions, each a real gap:
\begin{enumerate}[leftmargin=2em]
  \item \textbf{Color--valued.} The current $\sigma\sigma$ acts on the spin/charge
        qubits. A gluon's exponential must act on the \emph{color register} --- the
        $b_3$ harmonic sector carrying $[1,\omega,\omega^2]$ --- a different
        bundle from the spin fiber.
  \item \textbf{Non--abelian (root, not Cartan--diagonal).} A photon is the
        abelian direction; a gluon is built from the off--diagonal roots
        $E_{\pm\alpha}$, so its holonomy is \emph{path--ordered} and
        non--commuting and reproduces the $\Z_3/S_3$ color cycling around a loop
        (cf.\ $M$ in \eqref{eq:intertwine}). The neutral $q_{AB}=0$ core is, by
        contrast, abelian.
  \item \textbf{Geometry--sourced, on the right connection.} The existing
        inter--hole Wilson line (\code{wilson\_line}, \code{facet\_transport}) is
        the $\Spin(4)$ \emph{spin connection} --- the deficit--angle, gravity--like
        frame alignment (\code{dk\_spin\_wilson.py}). It is $K$--type (local,
        non--entangling) and lives on the \emph{frame} bundle, not the color one.
        A gluon needs a \emph{color} connection on the register, whose holonomy
        supplies the chamber point $c$. \textbf{Do not conflate the two}: aligning
        spinor frames can never entangle (Claim~\ref{cl:onlyA}); that is the floor,
        not the cure.
\end{enumerate}

\paragraph{Bare gluon vs.\ effective hyperfine --- the honest interpretation.}
Even granted (1)--(3), $\exp(i\sum_a c_a\,\sigma_a\!\otimes\!\sigma_a)$ most
closely resembles not the fundamental colored gluon field (a color--octet vector)
but the \emph{one--gluon--exchange color--magnetic hyperfine} interaction
\cite{drgg}
\begin{equation}
  H_{\text{hyp}} \;\propto\; \sum_{i<j}\,
     (\bm\lambda_i\!\cdot\!\bm\lambda_j)\,
     (\bm S_i\!\cdot\!\bm S_j)\,\big/(m_i m_j),
  \label{eq:hyp}
\end{equation}
with $\bm\lambda$ the Gell--Mann color generators. On a color singlet
$\expval{\bm\lambda_i\!\cdot\!\bm\lambda_j}$ is a fixed constant, so the spin
dependence is $\propto\sum_{i<j}\bm S_i\!\cdot\!\bm S_j$, which is exactly the
operator that splits the nucleon from the $\Delta$:
\begin{equation}
  \textstyle\sum_{i<j}\bm S_i\!\cdot\!\bm S_j
   = \halff\big(J^2-\tfrac94\big)
   = \begin{cases} -\tfrac34 & \text{proton } (J^2=\tfrac34),\\[2pt]
                   +\tfrac34 & \Delta\ (J^2=\tfrac{15}{4}).\end{cases}
  \label{eq:NDelta}
\end{equation}
The $\sigma_a\!\otimes\!\sigma_a$ generator of the chamber \emph{is} the
spin--spin coupling of \eqref{eq:hyp}. So the precise, defensible statement is:
the entangling exponential acts like the \emph{gluon--induced spin--spin
(hyperfine) interaction} --- the very term whose sign in \eqref{eq:NDelta}
\emph{is} the proton/$\Delta$ distinction. In an \emph{emergent} theory this is
the right target: an effective gluon, sourced by register holonomy, would be a
\emph{derivation} of the hyperfine interaction from geometry rather than a
postulate. Calling it ``a gluon'' is justified up to that ``effective/emergent''
qualifier; calling it the bare colored field is not.

% =============================================================================
\section{Two payoffs that are actually usable}
\label{sec:payoff}
% =============================================================================

\paragraph{(a) The chamber coordinate is the gauge--invariant observable we want.}
\code{composite\_proton\_spin\_findings.md} (\#485) makes GAUGE-- and
RELABEL--invariance the mandatory gates for any $J^2$. The Weyl--chamber point
$(c_1,c_2,c_3)$ --- equivalently the Makhlin local invariants \cite{makhlin} --- is
\emph{by construction} invariant under per--hole frame changes
$U\mapsto(k\otimes k')U(k''\otimes k''')$. So it is automatically the kind of
observable those gates demand. Concretely, the chamber content of a hole pair is
carried by the \emph{connected} two--hole correlator
\begin{equation}
  C_{ij} \eqdef \expval{\bm S_i\!\cdot\!\bm S_j}
              - \expval{\bm S_i}\!\cdot\!\expval{\bm S_j},
  \label{eq:Cij}
\end{equation}
which vanishes identically on a product (the floor: $C_{ij}=0$) and is precisely
what the per--hole extraction throws away. Since
\begin{equation}
  J^2 = \tfrac94 + 2\!\sum_{i<j}\expval{\bm S_i\!\cdot\!\bm S_j}
      = \tfrac94 + 2\!\sum_{i<j}\big(\expval{\bm S_i}\!\cdot\!\expval{\bm S_j}+C_{ij}\big),
\end{equation}
restoring the $C_{ij}$ restores $\tfrac34$. This suggests a readout that may be
\emph{lighter} than the full total--space operator of \code{spin\_readout.tex}:
read the three \emph{pairwise joint reduced states} $\rho_{ij}$ off the single
carried field and extract $C_{ij}$, rather than assembling one $16\times16$ operator
on the joint field. The quantum subsystem already stores two--body joints
(\code{jointAB}/\code{rhoAB} in \code{InteractionSimulation.cpp}); the machinery
exists, and only the proton readout currently replaces it with a \code{kron}.
\emph{Caveat:} this is a sibling of the total--space plan, not a free pass --- it
still requires $\rho_{ij}$ to be the faithful joint two--hole state, which is the
same hard kernel (the Whitney / K\"ahler--Atiyah fiber$\leftrightarrow$cells lift).

\paragraph{Empirical support already on file.}
\code{composite\_proton\_spin\_findings.md} reports that the joint--state prototype
(\code{joint\_*}: per--hole spins read from the \emph{single}
$\code{carriedRepresentative}([h_0,h_1,h_2],[1,\omega,\omega^2])$ rather than three
independent extractions) is genuinely different from the product, ``leans toward
the proton channel'' ($J^2$ drops $1.80\!\to\!1.64$ on the cached structure as the
$120^\circ$ color phases impose ${\sim}120^\circ$ relative spins), and can produce
a \emph{definite} $\Delta$ ($J^2=3.75$) where the product cannot. That is exactly
the signature of beginning to keep the $A$ factor: retaining the joint object moves
$J^2$ off the product corner toward $\tfrac34$. (It is not yet a clean observable
--- it inherits relabel--sensitivity from multi--hole
\code{carriedRepresentative} --- which is the open work, not a refutation.)

\paragraph{(b) Promote the transport intertwiner to the $A$ factor.}
$M$ in \eqref{eq:intertwine} already implements the color $\Z_3$ rotation between
holes; today it is used only as a \emph{diagnostic} (the intertwining residual
${\sim}4\times10^{-2}$). The proposal is to use $M$ (or the full color holonomy it
generates) as the \emph{source} of the entangling exponential in the state
reconstruction --- the discrete realization of the off--diagonal--root transport
of \S\ref{sec:weyl}. This is the concrete bridge from ``we have a color--cycling
transport'' to ``we have a gluon--like $A$ factor.''

% =============================================================================
\section{Honest caveats and where the lens does \emph{not} help}
\label{sec:caveats}
% =============================================================================

\begin{enumerate}[leftmargin=2em]
  \item \textbf{KAK is two--party; the proton is three.} There is no single
        three--party Cartan decomposition with one Weyl chamber; three--qubit
        entanglement is genuinely richer (the $\mathrm{GHZ}$ vs.\ $W$ SLOCC
        classes). The clean realization is therefore a \emph{pairwise / diquark
        sequence} --- entangle the like ($uu$) pair into the spin--triplet, then
        couple the $d$ --- which is exactly the two--step build already on record
        ($q\!+\!q\to\text{diquark}$, $\text{diquark}\!+\!q\to\text{proton}$, \#489).
        Good: it matches the build. Bad: ``one Cartan decomposition of the proton''
        is not a literal object; the payoff in \S\ref{sec:payoff}(a) is pairwise
        for this reason.
  \item \textbf{Do not set $c$ by hand.} Choosing the chamber point that forces
        $\tfrac34$ would be inserting the answer by fiat --- the same mistake that
        had \code{relax\_singlet} removed for placing $[1,\omega,\omega^2]$ by
        hand. The discipline: the $c_a$ must be \emph{sourced} from the emergent
        color holonomy (payoff (b)) and the relaxed geometry then \emph{checked}
        against the proton vs.\ $\Delta$ chamber point. The lens supplies a
        diagnosis and the right invariant; it does not license hand--tuning.
  \item \textbf{It does not remove the hard assembly --- a different axis.}
        Two orthogonal problems are easy to conflate, so name them. The
        fiber$\leftrightarrow$cells lift is \emph{vertical}: at one location,
        gather $\psi$'s degree--stratified components (the degree--$0,1,2,3,4$
        pieces sitting on different simplices) into the single $16$--dimensional DK
        fiber where $\Sigma_a$ is even defined --- the Whitney / K\"ahler--Atiyah
        interpolation, a \emph{local} operation building the fiber over a point.
        The Cartan / nonlocal content is \emph{horizontal}: across \emph{different}
        holes, combine those fibers while keeping the cross--correlations --- the
        $A$ factor. KAK is a purely horizontal tool; fiber$\leftrightarrow$cells is
        the vertical one, and solving the horizontal does not touch the vertical.
        Worse for any hope that it might: in the operator route
        $J^2=\sum_a(\sum_h P_h\Sigma_a P_h)^2$ the vertical is logically
        \emph{prior} --- one cannot even write $\Sigma_a P_h$ (the spin generator
        restricted to a hole) until $\Sigma_a$ is defined on the assembled field,
        so the entangling cross terms sit \emph{downstream} of the lift. In code,
        \code{emergent\_spinor} already performs a degree--$3$--only slice of the
        vertical lift (least--squares the tetrahedral trivector minors into the
        Clifford algebra), which is why $|\expval{\bm S}|=\halff$ per hole is
        robust; the gap is the \emph{all--degree, jointly--across--holes} version.
        Reading two holes independently gives $\rho_i\otimes\rho_j$, so $C_{ij}=0$
        by construction. The pairwise--$C_{ij}$ route therefore \emph{reduces} the
        burden (a two--hole joint lift instead of the full three--hole all--degree
        global lift) but does not remove it --- unless the escape hatch below holds.
  \item \textbf{Spin connection $\ne$ color connection.} The only inter--hole
        connection currently in the mesh is the $\Spin(4)$ spin connection
        (gravity--like, $K$--type). The gluon program presupposes a \emph{color}
        connection on the $b_3$ register that does not yet exist as a first--class
        object; \eqref{eq:intertwine}'s $M$ is the nearest seed.
  \item \textbf{A vielbein alone is the floor.} The user's phrase ``vielbein from
        one hole to the next'' should be read carefully: a vielbein is a local
        frame ($K$, the \code{cell\_frame}); a pure frame change can never
        entangle. The new physics is the $A$ factor \emph{sandwiched between} the
        vielbeine, not the vielbeine themselves. Thinking ``vielbein'' alone
        rebuilds exactly the product floor.
\end{enumerate}

\paragraph{The one escape hatch.}
There is exactly one way the horizontal (nonlocal) perspective could
\emph{sidestep} --- not solve --- the vertical assembly: read the spin content
from a \emph{holonomy} rather than from a reconstructed point--fiber. A Wilson
loop or the intertwiner $M$ of \eqref{eq:intertwine} is intrinsically a
\emph{path} object, computed from transport, never from gathering a $16$--component
column at a point (this is the \code{dk\_spin\_wilson.py} / \#477 thread, ``a
frame--independent, register--specific spin holonomy''). The hatch is open
\textbf{iff} the connected correlator $C_{ij}$ is a \emph{holonomy invariant} ---
expressible purely from the chamber coordinates of the inter--hole transport (the
color $\Z_3$ phases of $M$, the spin eigenphases of the loop), with no contraction
against endpoint spinors. If it is, $J^2$ can be read off transport data and the
point--fiber assembly is avoided; if it is not, a (reduced, two--hole) vertical
lift remains. Which of these holds is a concrete, decidable question, set as the
core test of \S\ref{sec:program}.

% =============================================================================
\section{A concrete, falsifiable next experiment}
\label{sec:program}
% =============================================================================

In increasing order of commitment:
\begin{enumerate}[leftmargin=2em]
  \item \textbf{Pairwise--$C_{ij}$ readout against the validated instrument.}
        Implement \eqref{eq:Cij}--style pairwise correlators on the joint reduced
        states and confirm, against the hand--fed clean states already in
        \code{test\_dk\_joint\_spin.py}, that they return
        $\tfrac34,\ \tfrac74,\ \tfrac{15}{4}$ for proton / product / $\Delta$.
        This is cheap and either reproduces the instrument in chamber coordinates
        or falsifies the framing immediately.
  \item \textbf{$C_{ij}$ two ways --- the escape--hatch test.} On one fixture,
        compute $C_{ij}$ \emph{(a)} from the reconstructed joint two--hole reduced
        state $\rho_{ij}$ (the vertical route) and \emph{(b)} from the
        $M$/Wilson--loop holonomy invariants alone (the horizontal route), and
        compare. Agreement certifies that $C_{ij}$ is a holonomy invariant --- the
        escape hatch is open and the point--fiber assembly is avoidable;
        disagreement quantifies exactly how much vertical lift the holonomy cannot
        replace. Either outcome is a decisive, reportable result.
  \item \textbf{Read $C_{ij}$ off the emergent joint field.} Apply the same
        readout to the \code{joint\_*} carried representative and check whether
        $C_{ij}\ne0$ and whether $J^2$ moves further toward $\tfrac34$ than the
        $1.80\!\to\!1.64$ already seen --- with the GAUGE/RELABEL gates enforced.
  \item \textbf{Source $c$ from the color holonomy.} Define the color connection
        whose discrete transport is $M$ of \eqref{eq:intertwine}, measure its
        loop holonomy, map it to a chamber point, and test whether the relaxed
        proton geometry lands on the proton corner and the $\Delta$ geometry on
        the $\Delta$ corner --- \emph{without} the chamber point being put in by
        hand. Success here is the operational meaning of ``the exponential acts
        like a gluon.''
\end{enumerate}

\paragraph{One sentence.}
The Cartan/Weyl lens does not by itself produce $\tfrac34$, but it names the
missing object (the abelian factor $A$), gives it a gauge--invariant coordinate
(the Weyl chamber / $C_{ij}$), identifies the color group's Weyl $S_3$ with the
very swaps that produce $\tfrac34$, and turns ``is it a gluon?'' into a testable
program: source that coordinate from an emergent color holonomy and check it lands
on the proton corner.

% =============================================================================
\begin{thebibliography}{9}
\addcontentsline{toc}{section}{References}

\bibitem{khaneja} N.~Khaneja and S.~J.~Glaser,
  \emph{Cartan decomposition of $\mathrm{SU}(2^n)$ and control of spin systems},
  Chem.\ Phys.\ \textbf{267}, 11 (2001).

\bibitem{kraus} B.~Kraus and J.~I.~Cirac,
  \emph{Optimal creation of entanglement using a two--qubit gate},
  Phys.\ Rev.\ A \textbf{63}, 062309 (2001).

\bibitem{makhlin} Yu.~Makhlin,
  \emph{Nonlocal properties of two--qubit gates and mixed states, and the
  optimization of quantum computations},
  Quantum Inf.\ Process.\ \textbf{1}, 243 (2002).

\bibitem{drgg} A.~De~R\'ujula, H.~Georgi, and S.~L.~Glashow,
  \emph{Hadron masses in a gauge theory},
  Phys.\ Rev.\ D \textbf{12}, 147 (1975) --- one--gluon--exchange color--magnetic
  hyperfine interaction and the $N$--$\Delta$ splitting.

\bibitem{helgason} S.~Helgason,
  \emph{Differential Geometry, Lie Groups, and Symmetric Spaces},
  Academic Press (1978) --- Cartan decomposition, Weyl chambers.

\bibitem{companion} Tessera theory notes,
  \emph{From Spin--\textonehalf{} to the Entangled Readout of the Proton Spin}
  (\code{spin\_readout.tex}, this directory); design notes
  \code{joint\_proton\_spin\_findings.md} (\#489),
  \code{composite\_proton\_spin\_findings.md} (\#485),
  \code{proton\_symmetric\_windows.tex} (\#398).

\end{thebibliography}

\end{document}
